\subsection{Laser Ultrasonic NDT Fundamentals}
\label{sec:lu-fundamentals}

Laser ultrasonic testing generates elastic waves through rapid thermal expansion at the material surface. Two generation regimes exist: thermoelastic ablation, where the laser pulse induces stress without material removal, and ablative generation, where plasma formation ejects surface material to produce stronger acoustic waves \cite{liu2022deep}. The resulting waveforms are broadband and contain dispersive Lamb wave modes, whose phase velocities vary with frequency and plate thickness. This wave character enables defect detection in thin-walled structures but complicates signal interpretation due to mode superposition \cite{zhang2020deep}.

Common inspection targets include aluminum alloys and carbon fiber reinforced polymer (CFRP) composites, where laser ultrasound detects delaminations, voids, and bondline defects without requiring couplant. However, acquired signals suffer from low signal-to-noise ratio (SNR) due to multiple noise sources: laser-induced plasma emission, ambient acoustic interference, and photodetector noise. Consequently, signal averaging is standard practice---repeated acquisitions are ensemble-averaged to suppress random noise and improve SNR \cite{liu2022deep}. This averaging approach increases inspection time linearly and may induce surface modification on sensitive materials, motivating research into single-shot denoising alternatives.
